\noindent\emph{证明.}
Banach 空间 $L^B\!\left(\frac1{p-1}\right)$ 都是 Serre~[7] 意义下的 $B(I)$ 型;
即它们具有标准正交基, 也就是属于其中的单项式

\begin{equation}\tag{4.6.9}
\pi^{U_0}X_0^{U_0}X_1^{U_1}\cdots X_{n+2}^{U_{n+2}}
\end{equation}

\noindent
按定义, 称其为标准正交基, 意即

\begin{equation}\tag{4.6.10}
\operatorname{ord}\!\left(
\sum_U C_U\pi^{U_0}X_0^{U_0}\cdots X_{n+2}^{U_{n+2}}
\right)
=\min_U\operatorname{ord}(C_U).
\end{equation}

设

\begin{equation}\tag{4.6.11}
R_B=\text{$L^B\!\left(\frac1{p-1}\right)$ 中由阶 $\geq0$ 的元素组成的 $\mathcal O$ 子模}.
\end{equation}

要证明引理, 只需证

\begin{equation}\tag{4.6.12}
jd<\operatorname{card}(B)
\quad\Longrightarrow\quad
(\alpha_B)^{\operatorname{ord}(q)}R_B\subset q^jR_B,
\end{equation}

或更强地, 只需证

\begin{equation}\tag{4.6.13}
jd<\operatorname{card}(B)
\quad\Longrightarrow\quad
\alpha_B(R_B)\subset p^jR_B.
\end{equation}

由 (4.5.22),

\begin{equation}\tag{4.6.14}
\alpha_B=\frac1p\,\tau^{-1}\circ\psi_p\circ H_B(X).
\end{equation}

由于

\begin{equation}\tag{4.6.15}
H(X)\in L\!\left(\frac1{p-1}\right),
\qquad \operatorname{ord}(H(X))\geq0,
\end{equation}

有

\begin{equation}\tag{4.6.16}
H_B(X)\,R_B\subset R_B
\end{equation}

且显然,

\begin{equation}\tag{4.6.17}
\tau^{-1}(R_B)=R_B.
\end{equation}

因此, (4.6.13) 由下式推出

\begin{equation}\tag{4.6.18}
jd<\operatorname{card}(B)
\quad\Longrightarrow\quad
\psi_p(R_B)\subset p^{j+1}R_B.
\end{equation}

逐个单项式验证 (4.6.18). 有

\begin{equation}\tag{4.6.19}
\begin{cases}
\psi_p\!\left(\pi^{U_0}X_0^{U_0}\cdots X_{n+2}^{U_{n+2}}\right)=0,
&\text{除非 }p\mid U_i\text{ 对每个 }i,\\[3pt]
\psi_p\!\left(
\pi^{pU_0}X_0^{pU_0}\cdots X_{n+2}^{pU_{n+2}}
\right)
=\pi^{(p-1)U_0}
\underbrace{\pi^{U_0}X_0^{U_0}\cdots X_{n+2}^{U_{n+2}}}_{\in R_B}.
\end{cases}
\end{equation}

$\pi^{U_0}X_0^{U_0}\cdots X_{n+2}^{U_{n+2}}$ 属于
$L^B\!\left(\frac1{p-1}\right)$ 的条件为

\begin{equation}\tag{4.6.20}
\begin{cases}
dU_0=\displaystyle\sum_{i\geq1}U_i,\\
U_i>0\Longleftrightarrow i\in B,\qquad i\geq1.
\end{cases}
\end{equation}

特别地,

\begin{equation}\tag{4.6.21}
dU_0\geq\operatorname{card}(B).
\end{equation}

故

\begin{equation}\tag{4.6.22}
\operatorname{ord}\!\left(\pi^{(p-1)U_0}\right)
=U_0\geq\frac{\operatorname{card}(B)}d>j,
\end{equation}

并且由于 $U_0\in\mathbb Z$,

\begin{equation}\tag{4.6.23}
\operatorname{ord}\!\left(\pi^{(p-1)U_0}\right)
=U_0\geq j+1.
\end{equation}

结合 (4.6.19), 这就证明了 (4.6.18).

